% ============================================================
%  SLOTMEIN THESIS — PART 2: System Architecture + Engine
% ============================================================
% This file is intended to be \input{} into the main thesis.tex
% Chapters 3 and 4 are defined here.
% ============================================================

% =============================================================
%         CHAPTER 3 — SYSTEM ARCHITECTURE
% =============================================================
\chapter{System Architecture and Design}

\section{Overview}

SlotMeIn is architected as a layered, modular RESTful backend application. The overarching design philosophy is \textit{separation of concerns}: each component of the system is responsible for a well-defined business capability, interacts with other components through stable interfaces, and can be modified or extended independently. This philosophy is operationalised through the application of Domain-Driven Design (DDD) principles \cite{evans_2003}, which organise code around the natural conceptual entities of the scheduling domain rather than around generic technical tiers.

The system is deployed as a cloud-hosted API accessible at \texttt{https://slotmein.vercel.app}, with a live interactive documentation interface auto-generated by FastAPI at the \texttt{/docs} endpoint. The PostgreSQL database is hosted on Supabase, providing a managed, scalable relational store with built-in connection pooling and Row-Level Security.

\section{High-Level Architecture}

At the highest level of abstraction, SlotMeIn consists of four layers:

\begin{enumerate}
    \item \textbf{Authentication Layer} — Handles user identity, JWT issuance and validation, and access control.
    \item \textbf{Domain Layer} — Implements the core business logic across six bounded contexts: \textit{Talents}, \textit{Shift Periods}, \textit{Shift Templates}, \textit{Constraints}, \textit{Schedule}, and \textit{Utilities}.
    \item \textbf{Data Layer} — Manages database connectivity, ORM model definitions, and session lifecycle.
    \item \textbf{Scheduling Engine} — A self-contained algorithmic subsystem nested within the Schedule domain that implements the constraint-satisfaction and scoring logic.
\end{enumerate}

Client applications communicate with all layers exclusively through the HTTP/REST interface. There are no shared in-process objects between the client and the server; all state flows through the API. This stateless design—mandated by REST's architectural constraints \cite{fielding_2000}—ensures that the system scales horizontally and that each request is independently interpretable.

\section{Technology Stack}

Table~\ref{tab:tech_stack} summarises the technology choices and the rationale for each.

\begin{table}[H]
\centering
\caption{SlotMeIn Technology Stack}
\label{tab:tech_stack}
\begin{tabular}{>{\bfseries}p{3.5cm} p{3.5cm} p{6cm}}
\toprule
Component & Technology & Rationale \\
\midrule
Web Framework & FastAPI 0.110+ & Native async, auto OpenAPI docs, strong Pydantic integration \\
Language & Python 3.11+ & Rich ML/data science ecosystem; concise syntax \\
ORM & SQLAlchemy 2.0 & Modern mapped-column API; type-safe queries \\
Database & PostgreSQL & ACID compliance; strong support for relational scheduling data \\
Auth & JWT + OAuth2 & Stateless; standard; FastAPI native support \\
Hosting & Supabase + Render & Managed Postgres; auto-scaling server deployment \\
Testing & pytest 9.0 & Mature fixture system; parametrised tests \\
Data Science & pandas, NumPy & Data preprocessing and simulation \\
Email & Resend & Transactional email for user invitations \\
\bottomrule
\end{tabular}
\end{table}

\section{Directory Structure and Domain Organisation}

The project follows a feature-based directory structure, placing all files related to a business capability under a single directory tree. This contrasts with layer-based structures (e.g., separating all ``models'' and all ``services'' into global directories) and is directly motivated by DDD's concept of bounded contexts \cite{evans_2003}.

\begin{lstlisting}[language=bash, style=pythonstyle, caption={Top-level project directory structure}, label={lst:dir}]
app/
  authentication/       # Auth bounded context
    routes.py           # Login / token endpoints
    users/              # User CRUD
    tokens/             # JWT generation and validation
  config/               # Environment variable loading
  core/                 # Business domain contexts
    talents/            # Employee profile management
    shift_period/       # Shift window definitions
    shift_template/     # Per-role slot templates
    constraints/
      constraint_rules/   # Global labour rules
      talent_constraints/ # Per-talent rule assignments
    schedule/           # Scheduling bounded context
      allocator/        # Core engine
        engine/
          generators.py   # Eligible talent generation
          validators.py   # Hard constraint checkers
          scheduler_scoring.py  # Scoring + round-robin
      staffing/         # Staffing level determination
      shifts/           # Shift slot construction
      talents/          # Talent data pipeline
      routes.py         # Schedule API endpoints
  database/
    models.py           # SQLAlchemy ORM definitions
    session.py          # DB session factory
main.py                 # Application entry point
\end{lstlisting}

Each domain context exposes its own \texttt{routes.py} (API endpoints), \texttt{service.py} (business logic), \texttt{schema.py} (Pydantic models for request/response validation), and \texttt{repo.py} (data access). This clean separation ensures that changes to the database schema, for example, do not propagate unexpectedly into API response shapes.

\section{Database Schema}

The relational database schema consists of seven tables, each mapped to a Python class via SQLAlchemy's Declarative Base. Figure~\ref{fig:erd} conceptually represents the entity relationships.

\begin{table}[H]
\centering
\caption{Database Table Summary}
\label{tab:db_tables}
\begin{tabular}{>{\bfseries}p{4cm} p{9cm}}
\toprule
Table & Purpose \\
\midrule
\texttt{talents} & Employee profiles: name, role, contract type, weekly hours, active status \\
\texttt{talent\_constraints} & Associates a talent with a constraint type (e.g., ``availability'') \\
\texttt{constraint\_rules} & Per-constraint rules specifying available days and permitted shifts \\
\texttt{shift\_periods} & Defines named shift windows (e.g., ``Morning'', ``Evening'') \\
\texttt{shift\_templates} & Per-role time slots within a period \\
\texttt{scheduled\_shifts} & Individual talent-to-shift assignments within a saved schedule \\
\texttt{schedules} & Week-level scheduling documents with status (\textit{draft}/\textit{final}) \\
\bottomrule
\end{tabular}
\end{table}

\subsection{Key Design Decisions}

\paragraph{Cascade Deletion.} All child records (e.g., constraints, scheduled shifts) are configured with \texttt{ON DELETE CASCADE}, ensuring referential integrity when a parent talent or schedule is removed.

\paragraph{Draft/Final Status.} The \texttt{schedules} table maintains a \texttt{status} field that transitions from \texttt{draft} to \texttt{final} upon manager approval. A uniqueness guard prevents multiple final schedules for the same week, protecting against accidental overwrites.

\paragraph{Separation of Periods and Templates.} The distinction between \texttt{shift\_periods} (named windows with a general start/end time) and \texttt{shift\_templates} (per-role time slots within a period) provides a two-level configurability: managers can adjust role-level slot timings without redefining the named shift period, and vice versa.

\section{Authentication and Security}

The authentication system follows the OAuth2 Password Flow, in which the client submits username and password to a \texttt{/login} endpoint and receives a signed JWT access token. All subsequent API requests include this token in the \texttt{Authorization: Bearer} header.

Tokens are validated using the \texttt{python-jose} library against a server-side secret key. Password storage uses \texttt{bcrypt} via \texttt{passlib}, ensuring that plaintext passwords are never persisted. The \texttt{get\_current\_user} dependency—injected into every protected endpoint via FastAPI's dependency injection system—ensures that unauthenticated requests are rejected at the framework level before any business logic executes.

User creation follows an invite-only model: a superuser sends an invitation email via the Resend transactional email service contains a registration link. This design is appropriate for the managed, closed-team context of a hotel or restaurant operation.

\section{API Endpoint Design}

Table~\ref{tab:endpoints} lists the key API endpoints exposed by the scheduling domain, which is the primary focus of this thesis.

\begin{table}[H]
\centering
\caption{Schedule Domain API Endpoints}
\label{tab:endpoints}
\begin{tabular}{p{2cm} p{5cm} p{6cm}}
\toprule
\textbf{Method} & \textbf{Path} & \textbf{Description} \\
\midrule
POST & \texttt{/generate} & Run the scheduling algorithm; return a preview (not persisted) \\
POST & \texttt{/commit} & Persist a previewed schedule as draft or final \\
GET  & \texttt{/} & List all saved schedules \\
GET  & \texttt{/\{schedule\_id\}} & Retrieve a specific schedule with all assignments \\
PATCH & \texttt{/\{schedule\_id\}/status} & Promote a draft to final \\
DELETE & \texttt{/\{schedule\_id\}} & Delete a schedule and all its assignments \\
POST & \texttt{/assignments/} & Manually add a single assignment \\
PATCH & \texttt{/assignments/\{id\}} & Edit an existing assignment \\
DELETE & \texttt{/assignments/\{id\}} & Remove an assignment \\
POST & \texttt{/validate\_assignment} & Check a proposed manual assignment for constraint violations \\
\bottomrule
\end{tabular}
\end{table}

The separation between \texttt{/generate} (preview) and \texttt{/commit} (persist) is a deliberate UX-oriented design decision. Managers can review the algorithm's recommendations before committing them to the database, and can make manual adjustments—validated against the constraint engine via \texttt{/validate\_assignment}—before finalising the schedule.

\section{Data Flow for Schedule Generation}

The schedule generation flow, triggered by \texttt{POST /generate}, proceeds through the following pipeline:

\begin{enumerate}
    \item \textbf{Week Expansion:} A \texttt{weekRange} object is constructed from the input \texttt{start\_date}, producing the seven calendar dates of the target week.
    \item \textbf{Shift Slot Construction:} The \texttt{ShiftSlotBuilder} queries \texttt{shift\_periods} and \texttt{shift\_templates} from the database and constructs a dictionary of \texttt{shiftSpecification} objects keyed by a unique shift instance identifier.
    \item \textbf{Talent Pipeline:} The \texttt{TalentService} loads all active talents via \texttt{TalentRepository}, processes their availability windows through \texttt{TalentPreprocessor}, and assembles them into \texttt{talentAvailability} objects via \texttt{TalentAssembler}.
    \item \textbf{Historical Context:} The prior week's \texttt{ScheduledShift} records are loaded and converted into \texttt{assignment} objects that seed the engine with cross-week continuity (e.g., detecting consecutive day streaks that began in the previous week).
    \item \textbf{Scheduling Engine:} The \texttt{ScheduleBuilder} orchestrates the allocation algorithm, producing a list of \texttt{assignment} objects.
    \item \textbf{Understaffing Analysis:} The \texttt{UnderstaffedShifts} component identifies any shifts where required headcount exceeds assigned headcount.
    \item \textbf{Response Serialisation:} Assignments are serialised to JSON and returned to the client alongside understaffed shift alerts.
\end{enumerate}

% =============================================================
%        CHAPTER 4 — THE SCHEDULING ENGINE
% =============================================================
\chapter{The Deterministic Scheduling Engine}

\section{Overview}

The scheduling engine is the intellectual core of SlotMeIn. It is a self-contained algorithmic subsystem that accepts a set of shift requirements and a set of talent availability objects, and produces a legally compliant, fairness-aware assignment plan. The engine is deterministic—given the same inputs, it produces the same outputs—which makes it auditable, testable, and explainable.

The engine's design is inspired by the constraint-satisfaction literature \cite{rossi_et_al_2006} and by practical scheduling system design guidelines \cite{pinedo_2016}. It implements a \textit{priority-based constructive heuristic}: shifts are processed in order of increasing flexibility (measured by the number of eligible candidates), and for each shift a scoring function ranks candidates before a round-robin tiebreaker selects the final assignee.

\section{Core Data Abstractions}

Before describing the algorithm, it is important to understand the data abstractions on which it operates.

\subsection{shiftSpecification}

A \texttt{shiftSpecification} represents a concrete, time-stamped shift slot to be filled. It carries the following attributes:
\begin{itemize}
    \item \texttt{template\_id}: Reference to the originating shift template.
    \item \texttt{start\_time}, \texttt{end\_time}: Full \texttt{datetime} objects (date + time).
    \item \texttt{shift\_name}: Named shift period (e.g., \textit{Morning}, \textit{Evening}).
    \item \texttt{role\_name}: The role required (e.g., \textit{bartender}, \textit{server}).
    \item \texttt{role\_count}: The number of employees of this role required for this slot.
\end{itemize}

\subsection{talentAvailability}

A \texttt{talentAvailability} object encapsulates everything the engine needs to know about a talent:
\begin{itemize}
    \item \texttt{talent\_id}: Unique identifier.
    \item \texttt{role}: The role enum value.
    \item \texttt{weeklyhours}: The talent's contractual weekly hours cap.
    \item \texttt{window}: A dictionary mapping \texttt{date} objects to a list of \texttt{(start\_dt, end\_dt)} availability spans.
    \item \texttt{shift\_name}: The set of shift names the talent is eligible for.
    \item \texttt{constraint}: A flag indicating whether the talent has active scheduling constraints.
\end{itemize}

\subsection{assignment}

An \texttt{assignment} records a confirmed allocation: a \texttt{talent\_id}, a \texttt{shift\_id}, and the full \texttt{shiftSpecification} for that allocation. The list of assignments grows as the engine processes each shift, and is also seeded with prior-week history to support cross-week constraint checking.

\section{Staffing Level Determination}

Before the engine assigns individuals to shifts, the required headcount for each role on each day must be determined. This is the responsibility of the \texttt{StaffingService}.

The service applies a configurable day-of-week staffing rule:

\begin{table}[H]
\centering
\caption{Staffing Level Rules by Day of Week}
\label{tab:staffing}
\begin{tabular}{lll}
\toprule
\textbf{Level} & \textbf{Days} & \textbf{Staffing Multiplier} \\
\midrule
Low    & Monday, Tuesday       & Minimum headcount \\
Medium & Wednesday, Thursday   & Standard headcount \\
High   & Friday, Saturday, Sunday & Peak headcount \\
\bottomrule
\end{tabular}
\end{table}

The headcount for each role at each level is defined in a STAFFING\_CONFIGURATION\_RULES dictionary. For example, the \textit{Server} role requires 2 employees on low-demand days, 3 on medium, and 4 on high-demand days. This configuration is currently deterministic but is designed to be replaceable by AI-predicted demand values in a future iteration (see Chapter~\ref{ch:ai_readiness}).

\section{The Allocation Algorithm}

\subsection{Stage 1: Eligibility Generation}

The \texttt{TalentAvailabilityService} generates a dictionary mapping each shift instance identifier to an ordered list of eligible talent IDs. Eligibility is determined by the \texttt{TalentGenerator}, which applies three filters:

\begin{enumerate}
    \item \textbf{Role match:} The talent's role must match the shift's required role.
    \item \textbf{Availability window:} The talent must have an availability span on the shift's date that fully covers the shift's start and end times (i.e., $\text{avail\_start} \leq \text{shift\_start}$ AND $\text{avail\_end} \geq \text{shift\_end}$).
    \item \textbf{Shift name eligibility:} The talent must be eligible for the named shift (e.g., a talent who only works morning shifts is not eligible for evening slots).
\end{enumerate}

Within the eligible list, constrained talents (those with active \texttt{TalentConstraint} records) are placed before unconstrained ones, implementing a soft priority for talent-specific scheduling agreements.

\subsection{Stage 2: Least-Flexible-First Ordering}

The engine sorts shifts in ascending order of the size of their eligibility list:

\begin{lstlisting}[caption={Least-flexible-first shift ordering}, label={lst:lffo}]
sorted_shifts = sorted(
    self.assignable_shifts.items(),
    key=lambda shift: len(eligibility.get(shift[0], []))
)
\end{lstlisting}

This implements the well-known \textit{fail-first} (or \textit{most-constrained variable}) heuristic from constraint programming \cite{rossi_et_al_2006}: by processing the hardest-to-fill shifts first, the algorithm reduces the risk of reaching an unresolvable dead end for critical shifts after easier ones have consumed the available talent pool.

\subsection{Stage 3: Scoring}

For each shift, the \texttt{computeScore} engine calculates a composite score for every eligible candidate. The score consists of three additive components:

\paragraph{Hours Balance Score.} This penalises talents who are already close to their weekly hour limit:
\[
S_{\text{hours}}(t) = H_{\text{weekly}}(t) - H_{\text{assigned}}(t)
\]
where $H_{\text{weekly}}(t)$ is the talent's contracted weekly hours and $H_{\text{assigned}}(t)$ is the hours already allocated in the current scheduling run. A talent with more remaining hours receives a higher score, promoting equitable workload distribution aligned with contractual entitlements.

\paragraph{Consecutive Streak Score.} This penalises talents who have recently worked many consecutive days, discouraging near-maximal streaks before legal limits are reached:

\begin{table}[H]
\centering
\caption{Consecutive Streak Penalty Tiers}
\label{tab:streak}
\begin{tabular}{cc}
\toprule
\textbf{Streak (days)} & \textbf{Penalty} \\
\midrule
0   & $0.0$ \\
1--2 & $-2.0$ \\
3--4 & $-5.0$ \\
5--6 & $-10.0$ \\
\bottomrule
\end{tabular}
\end{table}

\paragraph{Rest Gap Score.} If the talent worked the previous day and the gap between yesterday's end time and today's proposed start time is less than 11 hours, a penalty of $-5.0$ is applied:
\[
S_{\text{rest}}(t) = \begin{cases} -5.0 & \text{if } \Delta t_{\text{rest}} < 11 \text{ hours} \\ 0.0 & \text{otherwise} \end{cases}
\]

The composite score is:
\[
S(t) = S_{\text{hours}}(t) + S_{\text{streak}}(t) + S_{\text{rest}}(t)
\]

\subsection{Stage 4: Round-Robin Tiebreaking}

After scoring, the candidate with the highest $S(t)$ is selected. When multiple candidates achieve the maximum score, the \texttt{roundRobinPicker} selects among them using a per-role pointer:

\begin{lstlisting}[caption={Round-robin picker implementation}, label={lst:rr}]
def pickBestFit(self, role: str, candidates: list[int]) -> int | None:
    if not candidates:
        return None
    idx = self.pointers.get(role, 0) % len(candidates)
    chosen = candidates[idx]
    self.pointers[role] = self.pointers.get(role, 0) + 1
    return chosen
\end{lstlisting}

The pointer advances modulo the number of tied candidates, cycling through them in order. This ensures that over a full week, tied candidates receive approximately equal shift counts—a fairness guarantee analogous to proportional representation in social choice theory \cite{bouveret_lemaitre_2016}.

\subsection{Stage 5: Hard Constraint Validation}

Before an assignment is recorded, the selected candidate passes through a chain of four hard constraint validators. These validators implement the \texttt{abstractValidator} interface, allowing new constraints to be added without modifying the allocation loop.

\begin{table}[H]
\centering
\caption{Hard Constraint Validators}
\label{tab:validators}
\begin{tabular}{p{4.5cm} p{8.5cm}}
\toprule
\textbf{Validator} & \textbf{Rule Enforced} \\
\midrule
\texttt{maxHoursValidator} & Total weekly hours (including this shift, minus break duration) must not exceed the talent's contracted limit plus a 1-hour buffer. \\
\texttt{dailyAssignmentValidator} & A talent may not be assigned more than one shift per calendar day. \\
\texttt{restValidator} & At least 11 hours must elapse between the end of the previous day's shift and the start of the proposed shift. \\
\texttt{consecutiveValidator} & A talent may not work more than 6 consecutive days. \\
\bottomrule
\end{tabular}
\end{table}

If any validator returns \texttt{False}, the candidate is discarded and the algorithm proceeds to the next candidate in the ranked list. This separation of scoring (soft preference) and validation (hard constraint) mirrors the classical constraint programming distinction between objective function and feasibility constraints \cite{pinedo_2016}.

\subsection{Stage 6: Assignment Recording}

If a candidate passes all validators, a new \texttt{assignment} is created and appended to both the plan (the final output) and the working assignments list (the running state). The talent's workload counter is incremented by the net shift hours (gross hours minus break duration):

\begin{lstlisting}[caption={Workload tracking and state update}, label={lst:workload}]
shift_hours = (shift.end_time - shift.start_time)\
              .total_seconds() / 3600
shift_hours -= get_break_duration(shift.shift_name)
workload[best_fit] += shift_hours

for validator in validators:
    if hasattr(validator, "mark_assigned"):
        validator.mark_assigned(ctx)
\end{lstlisting}

The \texttt{mark\_assigned} call updates stateful validators (specifically \texttt{dailyAssignmentValidator}, which maintains a set of \texttt{(talent\_id, date)} pairs) to reflect the new assignment.

\section{Understaffing Detection}

After the allocation loop completes, the \texttt{UnderstaffedShifts} component compares required headcount (from \texttt{shiftSpecification.role\_count}) against actual assigned count for each shift instance. Any gap is reported as an understaffed shift alert, giving the manager actionable visibility into coverage gaps that require manual intervention.

\section{Manual Assignment Validation}

The \texttt{POST /validate\_assignment} endpoint allows managers to check a proposed manual assignment against all four hard constraint validators before committing it to the database. Violations are returned as human-readable messages:

\begin{itemize}
    \item \textit{``Exceeds weekly hours — would exceed their Nh contract limit.''}
    \item \textit{``Already scheduled — already has a shift on [date].''}
    \item \textit{``Insufficient rest — less than 11 hours since their previous shift.''}
    \item \textit{``Too many consecutive days — this would be their 7th consecutive working day.''}
\end{itemize}

Violations are informational rather than blocking—managers can override the engine's recommendations when operational circumstances demand it—but the audit trail ensures that constraint violations are always visible and deliberate.

\section{Design Critique and Future Improvements}

The current deterministic engine is robust, explainable, and legally compliant, but it has known limitations:

\begin{enumerate}
    \item \textbf{Greedy irrecoverability:} The constructive heuristic does not backtrack. If an early assignment consumes a talent needed for a later, more critical shift, the algorithm cannot undo that decision. A future version could implement limited backtracking or a local search post-processing step.
    \item \textbf{Static staffing levels:} The day-of-week demand tiers (low/medium/high) do not account for date-specific factors such as public holidays, local events, or seasonal patterns. This is the primary motivation for the AI predictive layer described in Chapter~\ref{ch:ai_readiness}.
    \item \textbf{Single shift per day constraint:} The current \texttt{dailyAssignmentValidator} prevents any talent from working more than one shift per day, even if their availability and hours would permit a split shift. This is intentionally conservative but may exclude valid scheduling options in certain operational contexts.
    \item \textbf{Scalability:} The $O(S \cdot T)$ scoring loop (over shifts $S$ and talents $T$) is efficient for typical restaurant-scale instances (tens of talents, hundreds of shifts per week) but would require optimisation for very large deployments.
\end{enumerate}

\end{document}
